\documentclass{rrxiv}
\rrxivid{rrxiv:2605.00007}
\rrxivversion{v1}
\rrxivprotocolversion{0.1.0}
\rrxivlicense{CC-BY-4.0}
\rrxivtopics{cs.DL,cs.CY}

\title{Retraction notices as first-class data}
\author{H. Andersen \and C. Park}
\date{2026-05-16}

\begin{document}
\maketitle

\begin{abstract}
Most preprint servers handle retractions as freeform text. We argue retractions should be structured artifacts linked to the specific claims they invalidate, with machine-readable rationale and an audit trail.
\end{abstract}

\section{Claims}
\begin{claim}[Claim 1]
\label{claim:c1}
Structured retraction notices reduce downstream citation contamination by an estimated 71\%.

\emph{Replication status: untested.}
\end{claim}

\begin{claim}[Claim 2]
\label{claim:c2}
A retraction-as-annotation model is implementable without changing the underlying paper schema.

\emph{Replication status: untested.}
\end{claim}

\section{Notes}
This document is a synthetic render of the canonical CIR for demonstration purposes. The full source, claim graph, and discussion live in the rrxiv reference instance. See the \texttt{rrxiv} protocol repository for the schema.

\end{document}
